\chapter{Les articles lus}
\label{chap:C-articles}

Cette annexe liste les articles des deux échantillons de la comparaison principale, avec leur
identifiant, leur titre et leur conférence. L'identifiant renvoie au fichier de lecture
\file{study/codes/<lecteur>/<id>.json} ou \file{study/codes\_v2/<lecteur>/<id>.json}, où chaque
décision s'appuie sur une citation et sa page. Les PDF ne sont pas dans le dépôt
(\cref{sec:absent}). Le script \file{study/fetch.py} les télécharge depuis les adresses notées dans les
fichiers d'échantillon.

\section{Les articles de 2026}
\label{sec:c-2026}

La population compte les 1\,252 articles d'ICLR 2026, d'ICML 2026 et de NeurIPS 2025 dont le titre
correspond à l'expression régulière du protocole, sur 17\,836 articles acceptés
(\cref{chap:03-sondage-methode}). Le préfixe de l'identifiant dit d'où vient l'article.
\begin{itemize}
  \item \texttt{LR}, \texttt{ML}, \texttt{NE} : le premier tirage (\texttt{random.Random(2018)}),
        17 articles d'ICLR 2026, puis ICML 2026 et NeurIPS 2025 (\file{study/final\_2026.json}).
  \item \texttt{MLR}, \texttt{NER} : les huit articles de remplacement d'ICML et de NeurIPS, tirés
        par \texttt{random.Random(2026)} dans la même conférence, quand l'article tiré n'avait pas de
        PDF librement accessible (\file{replacements\_2026.json}).
  \item \texttt{K} : la file de l'élargissement (A3), mélangée par
        \texttt{random.Random(2026)}, lue dans cet ordre après le tri sur le titre jusqu'à 36
        articles de RL classique. La file complète se trouve dans
        \file{classic\_queue\_2026.json}.
\end{itemize}
L'échantillon de RL classique de 2026 de la comparaison principale réunit les 14 articles
\texttt{classic} du premier tirage et les 36 de la file A3. Le \cref{tab:articles-2026} liste les
articles lus. Ces articles ont été lus surtout dans leur version arXiv, ceux de 2018 dans les actes.
Pour le seul critère Q5, l'écart 31 du sondage a relu les articles de 2026 dans la version des actes
et les \enquote{non} de 2018 dans leur dernière version arXiv (\file{pdf/<id>-proc.pdf} et
\file{pdf/<id>-arxiv.pdf}), avec l'empreinte de chaque fichier dans \file{study/pdfs.csv}.

\gentable{articles-2026}

\section{Les articles de 2018}
\label{sec:c-2018}

Notre lecture de 2018 compte 50 articles de RL d'ICML 2018, de NeurIPS 2018 et d'ICLR 2018, les trois
conférences que Pineau annonçait. Les 12 premiers forment la calibration du protocole, 6 d'ICML et
6 de NeurIPS (\file{study/calib\_2018.json}). Les 38 autres viennent de l'élargissement A2, tirés
dans l'ordre d'un mélange \texttt{random.Random(2018)} avec une répartition proportionnelle à la
taille de chaque conférence (écart 12). L'échantillon réunit 21 articles d'ICML 2018 (préfixe
\texttt{C1}), 20 de NeurIPS 2018 (\texttt{C2}) et 9 d'ICLR 2018 (\texttt{C3})
(\file{study/calib\_2018\_v2.json}). Les comparaisons de chaque critère comptent les articles où le
critère s'applique, si bien qu'un dénominateur peut descendre sous 50.

\gentable{articles-2018}
